\documentclass[11pt,letterpaper]{article}
% This source is intentionally restricted to pdfLaTeX.
\usepackage[T1]{fontenc}
\usepackage[utf8]{inputenc}
\usepackage{textcomp}
\usepackage{lmodern}
\usepackage[margin=0.82in,top=0.94in,bottom=0.91in,headheight=18pt]{geometry}
\usepackage{amsmath,amssymb}
\usepackage[table]{xcolor}
\definecolor{DeepNavy}{HTML}{14283F}
\definecolor{Teal}{HTML}{0A747C}
\definecolor{Gold}{HTML}{C78F36}
\definecolor{Slate}{HTML}{536477}
\definecolor{Pale}{HTML}{F4F7FA}
\definecolor{CodeEdge}{HTML}{D4DFE9}
\usepackage{microtype}
\usepackage{parskip}
\setlength{\parskip}{0.47em}
\usepackage{titlesec}
\titleformat{\section}{\normalfont\Large\bfseries\color{DeepNavy}}{\thesection}{0.62em}{}
\titlespacing*{\section}{0pt}{2.0ex plus .7ex minus .2ex}{0.75ex}
\titleformat{\subsection}{\normalfont\large\bfseries\color{Teal}}{\thesubsection}{0.6em}{}
\usepackage{longtable,booktabs,array,calc}
\renewcommand{\arraystretch}{1.22}
\AtBeginEnvironment{longtable}{\small}
\usepackage{etoolbox}
\makeatletter
\patchcmd\longtable{\par}{\if@noskipsec\mbox{}\fi\par}{}{}
\makeatother
\usepackage{listings}
\IfFileExists{upquote.sty}{\usepackage{upquote}}{}
\lstdefinestyle{suitecode}{
  basicstyle=\ttfamily\footnotesize,
  breaklines=true,breakatwhitespace=false,
  columns=fullflexible,keepspaces=true,showstringspaces=false,
  frame=single,rulecolor=\color{CodeEdge},
  backgroundcolor=\color{Pale},
  keywordstyle=\color{Teal}\bfseries,
  commentstyle=\color{Slate}\itshape,
  stringstyle=\color{DeepNavy},
  xleftmargin=5pt,xrightmargin=5pt,
  framexleftmargin=5pt,framexrightmargin=5pt,
  aboveskip=0.85em,belowskip=0.95em,
  tabsize=2
}
\lstset{style=suitecode}
\usepackage{enumitem}
\setlist[itemize]{leftmargin=1.6em,itemsep=2pt,topsep=3pt}
\setlist[enumerate]{leftmargin=1.8em,itemsep=2pt,topsep=3pt}
\providecommand{\tightlist}{\setlength{\itemsep}{0pt}\setlength{\parskip}{0pt}}
\newcommand{\passthrough}[1]{#1}
\usepackage{xurl}
\usepackage[colorlinks=true,linkcolor=Teal,urlcolor=Teal,citecolor=Teal,pdfborder={0 0 0}]{hyperref}
\urlstyle{same}
\usepackage{fancyhdr,lastpage}
\pagestyle{fancy}
\fancyhf{}
\fancyhead[L]{\footnotesize\sffamily\color{Slate} ECMA-262:2025 \textbullet{} Engineering Documentation}
\fancyhead[R]{\footnotesize\sffamily\color{Slate} 16th Edition}
\fancyfoot[L]{\footnotesize\sffamily\color{Slate} Independent technical companion}
\fancyfoot[R]{\footnotesize\sffamily\color{Slate}\thepage\ / \pageref{LastPage}}
\renewcommand{\headrulewidth}{0.35pt}
\renewcommand{\footrulewidth}{0pt}
\setlength{\emergencystretch}{3em}
\setcounter{tocdepth}{2}
\setcounter{secnumdepth}{2}
\begin{document}
\begin{titlepage}
\thispagestyle{empty}
\vspace*{2.8cm}
{\color{Gold}\rule{18mm}{3pt}}\par\vspace{18mm}
{\sffamily\Large\bfseries\color{Teal} SOFTWARE ENGINEERING DOCUMENTATION SUITE}\par
\vspace{8mm}
{\sffamily\huge\bfseries\color{DeepNavy} ECMAScript 2025 }\par
\vspace{4mm}
{\sffamily\fontsize{24}{29}\selectfont\bfseries\color{DeepNavy} Secure Coding Guide }\par
\vspace{10mm}
{\large\color{Slate} Language-aware threat modeling and controls }\par
\vspace{16mm}
{\color{Teal}\rule{\textwidth}{0.8pt}}\par
\vfill
{\sffamily\color{Slate}\large ECMA-262:2025 \textbullet{} 16th Edition \textbullet{} June 2025}\par
\vspace{3mm}
{\sffamily\color{Slate} October 2026 \textbullet{} Document 5 of 8}\par
\vspace{3mm}
{\small\color{Slate} Original engineering companion. The uploaded ECMA-262:2025 specification remains the authoritative source.}\par
\end{titlepage}
\pagenumbering{arabic}
\tableofcontents
\clearpage
\section{Security scope and applicability}\label{security-scope-and-applicability}

This security engineering guide uses ECMA-262:2025 to explain language mechanisms that affect attack surface, then adds defensive requirements from application-security practice. \textbf{ECMA-262 is a language specification, not a secure-coding standard or a source of application authorization requirements.} In this document \textbf{{[}SPEC{]}} = language behavior defined in the standard, \textbf{{[}RISK{]}} = possible exploit scenario depending on application use, and \textbf{{[}CONTROL{]}} = recommended security engineering measure. CWE, OWASP, GHAS, and CodeQL are external taxonomies/tools, not normative parts of ECMA-262.

\section{Threat model and trust boundaries}\label{threat-model-and-trust-boundaries}

\begin{longtable}[]{@{}
  >{\raggedright\arraybackslash}p{(\columnwidth - 6\tabcolsep) * \real{0.2500}}
  >{\raggedright\arraybackslash}p{(\columnwidth - 6\tabcolsep) * \real{0.2500}}
  >{\raggedright\arraybackslash}p{(\columnwidth - 6\tabcolsep) * \real{0.2500}}
  >{\raggedright\arraybackslash}p{(\columnwidth - 6\tabcolsep) * \real{0.2500}}@{}}
\toprule\noalign{}
\begin{minipage}[b]{\linewidth}\raggedright
Boundary
\end{minipage} & \begin{minipage}[b]{\linewidth}\raggedright
Untrusted control
\end{minipage} & \begin{minipage}[b]{\linewidth}\raggedright
Common attack class
\end{minipage} & \begin{minipage}[b]{\linewidth}\raggedright
Defensive principle
\end{minipage} \\
\midrule\noalign{}
\endhead
\bottomrule\noalign{}
\endlastfoot
HTTP or message payload to JS object & names, shapes, values & prototype manipulation, coercion gadgets & validate schemas and allowed keys \\
User string to regex/template & pattern fragments, delimiters & injection or resource exhaustion & literal escaping, budgets \\
String to interpreted program & expression/source code & code injection & prohibit dynamic evaluation \\
Object to privileged decision & getters/proxies & TOCTOU, side effects & authenticate context, snapshot trusted fields \\
Module graph to executable code & dependency integrity & malicious packages / substitution & pin, verify, isolate dependencies \\
Browser DOM & text/HTML content & cross-site scripting & context-appropriate sinks and CSP \\
Shared memory / cross-agent state & order and races & state desynchronization & explicit synchronization and isolation \\
\end{longtable}

\textbf{{[}CONTROL{]}} Every risk assessment must identify attacker capability, entry point, vulnerable dataflow, intended sensitive operation, and realistic control. A surprising language coercion without a reachable security boundary is not automatically an exploitable vulnerability.

\section{SEC-01: Dynamic code evaluation}\label{sec-01-dynamic-code-evaluation}

\textbf{{[}SPEC{]}} The \passthrough{\lstinline!eval!} and Function constructors invoke compilation or evaluation semantics (19--20), with host-defined controls also appearing in the spec. \textbf{{[}RISK{]}} Treating user text as JavaScript code can grant the supplied string the privileges of the execution context.

\begin{lstlisting}
// UNSAFE: an untrusted expression can execute JavaScript.
function computeUnsafe(input) {
  return Function(`return (${input})`)();
}

// SAFER: explicitly choose an operation; no parsing of source code.
const operations = new Map([
  ['add', (a, b) => a + b],
  ['multiply', (a, b) => a * b],
]);
function compute(action, a, b) {
  if (!operations.has(action)) throw new TypeError('unknown operation');
  if (![a, b].every(Number.isFinite)) throw new TypeError('numbers only');
  return operations.get(action)(a, b);
}
\end{lstlisting}

\textbf{{[}CONTROL{]}} Disallow \passthrough{\lstinline!eval!}, indirect eval, \passthrough{\lstinline!new Function!}, and string-to-code timer APIs where possible. Use purpose-built parsers for constrained mini-languages and reject unknown grammar. \textbf{Verification:} CodeQL/ESLint rules, call graph review, taint tests from user input to evaluation sinks. \emph{Relevant weaknesses:} CWE-94 and CWE-95 (external taxonomy).

\section{SEC-02: Prototype pollution and unsafe keys}\label{sec-02-prototype-pollution-and-unsafe-keys}

\textbf{{[}SPEC{]}} Objects carry prototypes and property descriptors, and assignment/property lookup may involve inherited properties (10, 20). \textbf{{[}RISK{]}} Unvalidated keys, object merges, or special property paths can alter object behavior or enable privilege confusion, depending on merge implementation and target shape.

\begin{lstlisting}
// UNSAFE: dynamic path handling with magic properties can be dangerous.
function assignUnsafe(target, key, value) { target[key] = value; }

// SAFER: allowlisted schema; no attacker-selected property paths.
function updateProfile(profile, changes) {
  const allowed = new Set(['displayName', 'locale']);
  for (const [key, value] of Object.entries(changes)) {
    if (!allowed.has(key) || typeof value !== 'string')
      throw new TypeError('invalid profile field');
    profile[key] = value;
  }
  return profile;
}
\end{lstlisting}

\textbf{{[}CONTROL{]}} Use \passthrough{\lstinline!Map!} or null-prototype dictionaries for arbitrary keys; reject dangerous path segments; employ schema validation before deep merge; use \passthrough{\lstinline!Object.hasOwn!} when own-property presence is intended; avoid assigning into global prototypes. \textbf{Verification:} payloads using \passthrough{\lstinline!\_\_proto\_\_!}, \passthrough{\lstinline!constructor!}, \passthrough{\lstinline!prototype!}, nested paths, inherited getters. \emph{Relevant weakness:} CWE-1321.

\section{SEC-03: Injection through regular-expression construction}\label{sec-03-injection-through-regular-expression-construction}

\textbf{{[}SPEC{]}} RegExp grammar and \passthrough{\lstinline!RegExp.escape!} are defined in 22.2. An escaped literal has distinct semantics from arbitrary regex source. \textbf{{[}RISK{]}} User-controlled regex fragments can change matching logic or cause excessive runtime, depending on surrounding pattern.

\begin{lstlisting}
// UNSAFE when needle should be literal:
function unsafeContains(text, needle) {
  return new RegExp(needle).test(text);
}
// Literal-safe composition for engines with RegExp.escape:
function literalContains(text, needle) {
  return new RegExp(RegExp.escape(needle), 'u').test(text);
}
\end{lstlisting}

\textbf{{[}CONTROL{]}} Prefer \passthrough{\lstinline!includes!} when only substring search is intended; otherwise escape literal fragments and constrain user-controlled regex complexity. \textbf{Verification:} metacharacters, slash, newline, Unicode, catastrophic backtracking patterns, timeout/resource budgets. \emph{Relevant weaknesses:} CWE-1333 and contextual injection categories.

\section{SEC-04: Coercion, getters, and authorization logic}\label{sec-04-coercion-getters-and-authorization-logic}

\textbf{{[}SPEC{]}} \passthrough{\lstinline!ToPrimitive!}, getters, proxy traps, loose equality, and object property reads can execute user code (7, 10, 13). \textbf{{[}RISK{]}} Mixing representations or treating object fields as benign values can cause inconsistent authorization decisions.

\begin{lstlisting}
// UNSAFE: coercion-based comparison may accept unexpected representations.
function canAccessUnsafe(userId, ownerId) { return userId == ownerId; }

// SAFER: validate canonical identifiers before comparison.
function canAccess(userId, ownerId) {
  if (typeof userId !== 'string' || typeof ownerId !== 'string')
    return false;
  return userId === ownerId;
}
\end{lstlisting}

\textbf{{[}CONTROL{]}} Use canonical string/integer formats and strict comparisons for principals, tenant IDs, roles, and resource keys. Fetch trusted identity data from a validated source; revalidate where access decisions cross async boundaries. \textbf{Verification:} 0 versus \passthrough{\lstinline!"0"!}, empty string, \passthrough{\lstinline!null!}, symbols, BigInt, objects with malicious \passthrough{\lstinline!valueOf!} or getters. \emph{Relevant areas:} CWE-863, CWE-20, CWE-697, according to exact flaw.

\section{SEC-05: Asynchrony, race conditions, and resource lifetime}\label{sec-05-asynchrony-race-conditions-and-resource-lifetime}

\textbf{{[}SPEC{]}} Promise jobs, \passthrough{\lstinline!await!}, iterators, and shared-memory atomics have defined semantics (9, 25, 27, 29). \textbf{{[}RISK{]}} An application may check permission before an \passthrough{\lstinline!await!}, then act after authorization or resource state changes, or may conflate promise settlement with cancellation.

\begin{lstlisting}
// Potential business-logic race: privilege may change while awaiting.
async function transferUnsafe(actor, account) {
  if (!actor.canTransfer) throw new Error('forbidden');
  await reserveFunds(account);
  await executeTransfer(account); // no fresh authorization / transaction
}
\end{lstlisting}

\textbf{{[}CONTROL{]}} Enforce authorization near the commit point; use transactional consistency/locks when necessary; use explicit host cancellation (\passthrough{\lstinline!AbortSignal!} where supported), cleanup in \passthrough{\lstinline!finally!}, and iterator \passthrough{\lstinline!return!}/resource-close behavior. \passthrough{\lstinline!Promise.race!} alone does not abort I/O. \textbf{Verification:} force interleavings, deliberately reject promises, simulate cancellation and shutdown, exercise shared-memory data races. \emph{Relevant weaknesses:} CWE-367 and CWE-362 when supported by exploit context.

\section{SEC-06: Deserialization and structured input}\label{sec-06-deserialization-and-structured-input}

\textbf{{[}SPEC{]}} JSON parsing and serialization are part of clause 25.5; a parsed JSON object is still just data and does not guarantee domain validity. \textbf{{[}RISK{]}} Unvalidated nested objects can compromise downstream assumptions, including merge behavior and unsafe template rendering.

\begin{lstlisting}
export function parseRequest(text) {
  const value = JSON.parse(text);
  if (value === null || Array.isArray(value) || typeof value !== 'object')
    throw new TypeError('expected record');
  if (!Object.hasOwn(value, 'action') ||
      (value.action !== 'view' && value.action !== 'edit'))
    throw new TypeError('invalid action');
  return { action: value.action };
}
\end{lstlisting}

\textbf{{[}CONTROL{]}} Validate types, allowed fields, nesting/size limits, and semantic constraints; never infer authentication from parsed fields; keep data separate from HTML/script/code interpreters. \textbf{Verification:} oversized/nested JSON, duplicate keys policy, unexpected arrays, Unicode control characters, property names that collide with prototypes.

\section{SEC-07: Supply-chain and host isolation}\label{sec-07-supply-chain-and-host-isolation}

\textbf{{[}HOST{]}} ECMA-262 does not define npm packaging, browser Content Security Policy, network fetch behavior, URL resolution, or Node.js process permissions. \textbf{{[}RISK{]}} Importing code crosses a dependency-trust boundary and can run module initialization code with host-granted privileges.

\textbf{{[}CONTROL{]}} Pin dependencies with lockfiles, validate package provenance, restrict workflow tokens, review transitive dependencies, enable software composition analysis, isolate untrusted code, and use host-specific sandbox/permissions. Test JSON import attribute support per runtime. Document any \passthrough{\lstinline!import()!} paths derived from user input and enforce a constrained map of allowed modules.

\section{AppSec detection and triage matrix}\label{appsec-detection-and-triage-matrix}

\begin{longtable}[]{@{}
  >{\raggedright\arraybackslash}p{(\columnwidth - 6\tabcolsep) * \real{0.2500}}
  >{\raggedright\arraybackslash}p{(\columnwidth - 6\tabcolsep) * \real{0.2500}}
  >{\raggedright\arraybackslash}p{(\columnwidth - 6\tabcolsep) * \real{0.2500}}
  >{\raggedright\arraybackslash}p{(\columnwidth - 6\tabcolsep) * \real{0.2500}}@{}}
\toprule\noalign{}
\begin{minipage}[b]{\linewidth}\raggedright
Detection signal
\end{minipage} & \begin{minipage}[b]{\linewidth}\raggedright
Confirm by
\end{minipage} & \begin{minipage}[b]{\linewidth}\raggedright
Preferred disposition
\end{minipage} & \begin{minipage}[b]{\linewidth}\raggedright
Evidence
\end{minipage} \\
\midrule\noalign{}
\endhead
\bottomrule\noalign{}
\endlastfoot
Dynamic code construction & trace data from untrusted source to eval sink & remediate if attacker controlled & dataflow + executable PoC \\
User-controlled object keys & inspect merge implementation and target prototype & remediate exploitable path & before/after prototype invariants \\
Dynamic regex & determine if input is literal or pattern & escape or constrain & payload and resource-cost test \\
Loose equality in access checks & verify identifier validation and authorization path & use canonical strict comparison & privilege escalation test \\
Async check/use gap & identify mutable state across await & transaction/recheck & forced race schedule \\
Risky dependency/module import & inspect source, integrity, permissions & pin/remove/isolate & provenance and SBOM evidence \\
\end{longtable}

\textbf{{[}CONTROL{]}} Findings should identify reachable entry point, security impact, prerequisite conditions, mitigating controls, and remediation acceptance tests. A syntax property of ECMAScript is not itself a finding. Assign CWE only when the specific weakness is demonstrated; CVE/GHSA IDs refer to affected products/packages, not generic language features.

\section{Secure code review procedure}\label{secure-code-review-procedure}

\begin{enumerate}
\def\labelenumi{\arabic{enumi}.}
\tightlist
\item
  \textbf{Inventory surfaces:} input parsers, JSON, templating, dynamic property paths, module loading, eval-like constructs, regex, serialization.
\item
  \textbf{Identify privileged sinks:} authorization, credentials, filesystem, network, DOM HTML, code execution, data mutations.
\item
  \textbf{Trace flows:} source, transformations, coercion, implicit getter/proxy invocation, sanitizer/validator, sink.
\item
  \textbf{Assess execution:} async boundaries, module initialization, cleanup paths, cross-realm assumptions, shared memory.
\item
  \textbf{Reproduce:} write the smallest failing test with realistic attacker control; avoid purely theoretical alerts.
\item
  \textbf{Remediate:} implement allowlists/typed contracts/transactions and verify both positive and negative tests.
\item
  \textbf{Prevent recurrence:} regression tests, security-focused static analysis, dependency review, least-privilege CI.
\end{enumerate}

\section{Security verification exercises}\label{security-verification-exercises}

\textbf{Exercise A -- key contamination.} Write a vulnerable deep-merge utility in a disposable process, construct a nested \passthrough{\lstinline!constructor.prototype!} payload, then replace merge with explicit schema projection. Prove that Object.prototype is unchanged; do not execute against real applications.

\textbf{Exercise B -- regex hardening.} Compare dynamic \passthrough{\lstinline!new RegExp(input)!} with \passthrough{\lstinline!RegExp.escape(input)!} and \passthrough{\lstinline!includes!}. Test literal \passthrough{\lstinline!a+b!}, \passthrough{\lstinline![!}, \passthrough{\lstinline!\\n!}, and Unicode input. Analyze remaining performance risks.

\textbf{Exercise C -- authorization.} Implement a secure tenant-resource lookup with canonical ID types, strict equality, and a fresh commit-time authorization check. Deliberately force a role change during awaited I/O.

\textbf{Exercise D -- agent-synchronization.} Build a host-managed shared-memory counter with \passthrough{\lstinline!Atomics.add!}; compare with non-atomic read/modify/write under parallel workers and record observed outcomes without treating one test as a proof.

\section{Audit-ready security checklist}\label{audit-ready-security-checklist}

\begin{itemize}
\tightlist
\item[$\square$]
  Reject or isolate arbitrary source evaluation; identify all dynamic-code sinks.
\item[$\square$]
  Validate unknown JSON/property shapes before merging or authorizing.
\item[$\square$]
  Use canonical IDs and strict comparisons for security principals.
\item[$\square$]
  Audit prototype keys, getters, proxy traps, and inherited properties.
\item[$\square$]
  Escape regex literals and bound computation; use direct string APIs when possible.
\item[$\square$]
  Recheck sensitive state at async commit points and synchronize shared memory.
\item[$\square$]
  Constrain dynamic imports and verify dependencies and runtime privileges.
\item[$\square$]
  Attach proof, boundary, CWE rationale, tests, and owner to significant findings.
\end{itemize}

\section{References and scope notes}\label{references-and-scope-notes}

\textbf{Language source:} ECMA-262:2025 (uploaded PDF), clauses 2, 6--7, 9--10, 13--16, 19--20, 22, 25, 27, 29. \textbf{External nonnormative references:} CWE https://cwe.mitre.org/ ; OWASP https://owasp.org/ ; CodeQL https://codeql.github.com/docs/ ; ESLint https://eslint.org/docs/latest/rules/no-eval and https://eslint.org/docs/latest/rules/no-implied-eval . Security controls listed here are engineering recommendations, not statements of ECMA-262 conformance.

\end{document}
